\documentclass[12pt,a4paper]{article}

\newcommand{\ZVTetelsorDataOnly}{}
\newcommand{\ZVTetelOne}{Adatok, adattípusok, adatműveletek és adatstruktúrák. Szám adattípus. Számrendszerek, konverziók. A logikai, halmaz, karakter, sztring absztrakt adattípusok és realizációjuk. A tömb (vektor, mátrix), rekord absztrakt adattípusok.}
\newcommand{\ZVTetelTwo}{Az algoritmus. Iteratív és rekurzív algoritmus. A számítógépes memória. Adat és program. Verem és procedúra. Az algoritmus lejegyzése. A folyamatábra és a pszeudokód. Elemi algoritmusok.}
\newcommand{\ZVTetelThree}{Strukturált programozás. Programgráf, valódi program, vezérlőgráf lebontása, strukturált program és annak formulája. Strukturált programgráf kialakítása, struktogram. Ciklikus bonyolultság és egyéb bonyolultsági tételek.}
\newcommand{\ZVTetelFour}{Számelméleti algoritmusok. Legnagyobb közös osztó, euklideszi és kibővített euklideszi algoritmus, lineáris kongruencia egyenletek. Multiplikatív inverz, moduláris hatványozás, Fermat prímteszt. RSA.}
\newcommand{\ZVTetelFive}{Rendezések: Beszúró rendezés. Az ,,Oszd meg és uralkodj!'' elv. Összefésülő rendezés. Gyors rendezés. Buborék rendezés. Shell rendezés. Minimum kiválasztásos rendezés. Négyzetes rendezés. Lineáris idejű rendezések: leszámláló rendezés, számjegyes rendezés. Időelemzéseik. Az összehasonlító rendezések időtétele.}
\newcommand{\ZVTetelSix}{Gráfalgoritmusok. Szélességi keresés. Mélységi keresés. Topologikus rendezés. Optimumfeladatok fákon (minimális feszítőfák, Kruskal- és Prim-algoritmus). Legrövidebb utak meghatározása (Dijkstra-algoritmus, Bellman--Ford-algoritmus, Floyd--Warshall-algoritmus).}
\newcommand{\ZVTetelSeven}{A párhuzamos algoritmusok alapfogalmai, PRAM modell, hatékonysági mértékek (munka, költség, gyorsítás, hatékonyság). Párhuzamos végrehajtás ábrázolási módjai. Komplexitás becslése, párhuzamos komplexitás. Amdahl törvénye. Pipeline párhuzamosítás, Lost update probléma bemutatása és megoldása.}
\newcommand{\ZVTetelEight}{Aggregáció, redukció, prefixszámítás párhuzamosítási lehetőségei. \texttt{CREW\_PREFIX}, \texttt{EREW\_PREFIX} és \texttt{OPTIMAL\_PREFIX} algoritmusok bemutatása. Összefésülő rendezés párhuzamosítási módja.}
\newcommand{\ZVTetelNine}{A busz, gyűrű, rács, tórusz, hiperkocka és fa topológiák bemutatása, csomópontok közötti távolságok jellemzése. A CSP, Publish--Subscribe és az Aktor modell. Az aszinkron programozás elemei: async, await kulcsszavak használata, coroutine-ok.}
\newcommand{\ZVTetelTen}{A Turing-gép fogalma, működése. A RAM-gép. Boole-függvények és logikai hálózatok.}
\newcommand{\ZVTetelEleven}{Algoritmikus eldönthetőség. Church-tézis. Rekurzív és rekurzívan felsorolható nyelvek, rekurzív illetve parciálisan rekurzív függvények. Nevezetes nyelvek (Az R, Re, coR, coRE nyelvosztályok és ezek kapcsolata) és bonyolultságuk. Algoritmikusan eldönthetetlen problémák. Polinomiális idejű algoritmusok.}
\newcommand{\ZVTetelTwelve}{Nemdeterminisztikus Turing-gépek, az NP és a coNP nyelvosztály. Példák NP-beli nyelvekre. A tanú-tétel. Nemdeterminisztikus algoritmusok bonyolultsága. NP-teljesség, Cook-tétel. Néhány NP-teljes probléma, Cook--Levin-tétel.}
\newcommand{\ZVTetelThirteen}{A Java programozási nyelv története, alapvető jellemzői. Egy Java program szerkezete (csomagok és fordítási egységek). Az objektumorientált programtervezés alapelveinek felsorolása. Objektumok tárolása és élettartama. Szemétgyűjtő mechanizmus. Kivételkezelés.}
\newcommand{\ZVTetelFourteen}{Az egységbezárás és az információrejtés objektumorientált alapelvek megvalósítása. Osztálydefiníció és példányosítás. Osztályszintű és példányszintű tagok, konstansok. Hozzáférési kategóriák és használatuk. Hivatkozás az osztály tagjaira.}
\newcommand{\ZVTetelFifteen}{Az öröklődés és a polimorfizmus objektumorientált alapelvek megvalósítása. Öröklődés szabályai, metódusok túlterhelése és felüldefiniálása. A referencia statikus és dinamikus típusa. Absztrakt osztály és interfész szerepe, összehasonlítása.}
\newcommand{\ZVTetelSixteen}{Az operációs rendszerbeli folyamat (processz) fogalma. Folyamatkontextus. A fonál (szál) fogalma. A processzállapotok, állapotváltások, processz futási módok. Taszk- és fonálállapotok, állapotátmenetek.}
\newcommand{\ZVTetelSeventeen}{A memóriamenedzselés feladatai. Memória mint erőforrás. Címleképzés fajtái. Lapozós virtuális memóriamenedzselés működése. Allokálás, nyilvántartás, címleképzés. Laphiba fogalma, kezelése, kilapozási stratégiák. Előnyök, hátrányok.}
\newcommand{\ZVTetelEighteen}{Fájlrendszer-megvalósítási feladatok. Jegyzékszerkezetek. Szabadblokk-menedzselési lehetőségek. Fájlattribútum-rögzítési lehetőségek (ezen belül a fájltestet képező blokkok rögzítési lehetőségei).}
\newcommand{\ZVTetelNineteen}{Relációs adatmodell, relációs struktúra és integritási feltételek. ER modell elemei jelentése és jelölése. Az ER modell konverziója relációs modellre. Relációs adatmodell műveleti része, relációs algebra.}
\newcommand{\ZVTetelTwenty}{Az SQL szabvány relációs adatbázis kezelő nyelv bemutatása, a DDL (objektum létrehozás, megszüntetés és szerkezet módosítás), DML (rekord felvitel, törlés és módosítás), DCL és a SELECT utasítások áttekintése és használata. Beágyazott SELECT használata. A relációs algebra műveleteinek megvalósítása SQL-ben.}
\newcommand{\ZVTetelTwentyOne}{Adatkezelés és adatbáziskezelés alapfogalmai, fileszervezési módszerek, B-fa index; adatbázis architektúra; A normalizálás szerepe a relációs adatbázis kezelésben. Redundancia oka és a megszüntetés lépései. Normálformák. Dekompozíció célja és szabályai, tételei.}
\newcommand{\ZVTetelTwentyTwo}{Számítógép-hálózatokhoz kötődő alapfogalmak és az ISO--OSI hivatkozási modell: topológia és méret szerinti osztályozásuk. Kapcsolástechnika szerinti osztályozásuk (vonalkapcsolás, üzenetkapcsolás, csomagkapcsolás és virtuális vonalkapcsolás). Az ISO--OSI hivatkozási modell szerkezete, rétegei és azok főbb funkciói.}
\newcommand{\ZVTetelTwentyThree}{A számítógép-hálózatok ISO--OSI hivatkozási modell adatkapcsolati rétegének közeghozzáférési alrétege: A főbb csatorna-megosztási osztályok (statikus-dinamikus, dinamikus versengő-determinisztikus) bemutatása és azok összevetése. Az IEEE 802.3 és az Ethernet (CSMA/CD) keretformátum, MAC címek, támogatott közegek és sebességek bemutatása, működés duplex csatorna esetén. Az IEEE 802.11 WLAN általános felépítése, az Access Point (AP) főbb funkciói. Az alkalmazott közeghozzáférési módszer (CSMA/CA), Distributed Coordination Function (DCF) és Point Coordination Function (PCF) üzemmódok.}
\newcommand{\ZVTetelTwentyFour}{A TCP/IP protokollszövet és az Internet. Az Internet hivatkozási modell (DoD) és az ISO--OSI hivatkozási modell összevetése. A TCP/IP protokollszövet főbb részei (ARP, RARP, IP, ICMP, TCP, UDP) és azok funkciói. Az internetes címzés és címosztályok: IPv4-címosztályok, maszk, subnet, supernet, osztály nélküli címzés (CIDR -- Classless Inter-Domain Routing) és változó alhálózat-méretek (VLSM -- Variable Length Subnet Mask), címek kiosztása, lokális címek és címfordítás (NAT). IPv6-címek, IPv6-címtípusok (unicast, multicast, anycast; link local, unique local, aggregatable global unicast address).}
\newcommand{\ZVTetelTwentyFive}{Szoftvertechnológia és a szoftverfolyamat fogalma. A szoftverfolyamat-fázisok részletes bemutatása: szoftverspecifikáció, tervezés, implementáció, validáció és evolúció. Az evolúciós modell ismertetése.}
\newcommand{\ZVTetelTwentySix}{Szoftverkövetelmény fogalma és osztályozásai. Követelmények általános problémái. Követelményfeltárás és -elemzés folyamata. Vezérlési stílusok.}
\newcommand{\ZVTetelTwentySeven}{UML diagramok bemutatása: Use case és osztálydiagram. A szekvenciadiagram. Komponensdiagram.}

\newcommand{\ZVTetelKiiras}[1]{%
  \ifcase#1\relax
  \or \ZVTetelOne
  \or \ZVTetelTwo
  \or \ZVTetelThree
  \or \ZVTetelFour
  \or \ZVTetelFive
  \or \ZVTetelSix
  \or \ZVTetelSeven
  \or \ZVTetelEight
  \or \ZVTetelNine
  \or \ZVTetelTen
  \or \ZVTetelEleven
  \or \ZVTetelTwelve
  \or \ZVTetelThirteen
  \or \ZVTetelFourteen
  \or \ZVTetelFifteen
  \or \ZVTetelSixteen
  \or \ZVTetelSeventeen
  \or \ZVTetelEighteen
  \or \ZVTetelNineteen
  \or \ZVTetelTwenty
  \or \ZVTetelTwentyOne
  \or \ZVTetelTwentyTwo
  \or \ZVTetelTwentyThree
  \or \ZVTetelTwentyFour
  \or \ZVTetelTwentyFive
  \or \ZVTetelTwentySix
  \or \ZVTetelTwentySeven
  \fi
}

\newcommand{\ZVTetelsorList}{%
\begin{enumerate}[leftmargin=7mm,label=\arabic*.,itemsep=2pt,topsep=2pt,parsep=0pt]
  \item \ZVTetelOne
  \item \ZVTetelTwo
  \item \ZVTetelThree
  \item \ZVTetelFour
  \item \ZVTetelFive
  \item \ZVTetelSix
  \item \ZVTetelSeven
  \item \ZVTetelEight
  \item \ZVTetelNine
  \item \ZVTetelTen
  \item \ZVTetelEleven
  \item \ZVTetelTwelve
  \item \ZVTetelThirteen
  \item \ZVTetelFourteen
  \item \ZVTetelFifteen
  \item \ZVTetelSixteen
  \item \ZVTetelSeventeen
  \item \ZVTetelEighteen
  \item \ZVTetelNineteen
  \item \ZVTetelTwenty
  \item \ZVTetelTwentyOne
  \item \ZVTetelTwentyTwo
  \item \ZVTetelTwentyThree
  \item \ZVTetelTwentyFour
  \item \ZVTetelTwentyFive
  \item \ZVTetelTwentySix
  \item \ZVTetelTwentySeven
\end{enumerate}%
}


\newcommand{\tetelsorrow}[2]{\textbf{#1.} & #2\tabularnewline}

\newcommand{\tetelsorpageone}{%
\begin{tabularx}{\textwidth}{|>{\raggedleft\bfseries}p{8mm}@{\hspace{5mm}}>{\raggedright\arraybackslash}X|}
\hline
\tetelsorrow{1}{\ZVTetelOne}
\tetelsorrow{2}{\ZVTetelTwo}
\tetelsorrow{3}{\ZVTetelThree}
\hline
\tetelsorrow{4}{\ZVTetelFour}
\tetelsorrow{5}{\ZVTetelFive}
\tetelsorrow{6}{\ZVTetelSix}
\hline
\tetelsorrow{7}{\ZVTetelSeven}
\tetelsorrow{8}{\ZVTetelEight}
\tetelsorrow{9}{\ZVTetelNine}
\hline
\tetelsorrow{10}{\ZVTetelTen}
\tetelsorrow{11}{\ZVTetelEleven}
\tetelsorrow{12}{\ZVTetelTwelve}
\hline
\end{tabularx}%
}

\newcommand{\tetelsorpagetwo}{%
\begin{tabularx}{\textwidth}{|>{\raggedleft\bfseries}p{8mm}@{\hspace{5mm}}>{\raggedright\arraybackslash}X|}
\hline
\tetelsorrow{13}{\ZVTetelThirteen}
\tetelsorrow{14}{\ZVTetelFourteen}
\tetelsorrow{15}{\ZVTetelFifteen}
\hline
\tetelsorrow{16}{\ZVTetelSixteen}
\tetelsorrow{17}{\ZVTetelSeventeen}
\tetelsorrow{18}{\ZVTetelEighteen}
\hline
\tetelsorrow{19}{\ZVTetelNineteen}
\tetelsorrow{20}{\ZVTetelTwenty}
\tetelsorrow{21}{\ZVTetelTwentyOne}
\hline
\tetelsorrow{22}{\ZVTetelTwentyTwo}
\tetelsorrow{23}{\ZVTetelTwentyThree}
\tetelsorrow{24}{\ZVTetelTwentyFour}
\hline
\tetelsorrow{25}{\ZVTetelTwentyFive}
\tetelsorrow{26}{\ZVTetelTwentySix}
\tetelsorrow{27}{\ZVTetelTwentySeven}
\hline
\end{tabularx}%
}

\ifdefined\ZVTetelsorDataOnly
\expandafter\endinput
\fi

\documentclass[10pt,a4paper]{article}
\newcommand{\ZVHEADERLEFT}{Záróvizsga tételek}
\newcommand{\ZVHEADERRIGHT}{Programtervező informatikus BSc}
\usepackage{fontspec}
\usepackage{geometry}
\usepackage{microtype}
\usepackage{xcolor}
\usepackage{amsmath,amssymb,mathtools}
\usepackage{booktabs,longtable,tabularx,array,multirow}
\usepackage{enumitem}
\usepackage{hyperref}
\usepackage{fancyhdr}
\usepackage{titlesec}
\usepackage{listings}
\usepackage[most]{tcolorbox}
\usepackage{tikz}
\usetikzlibrary{positioning,shapes.geometric,shapes.misc,arrows.meta,decorations.pathreplacing}
\usepackage{caption}

\setlength{\headheight}{16pt}
\emergencystretch=2em
\renewcommand{\contentsname}{Tartalom}

% Fontbeallitasok: ha a DejaVu fontok nincsenek telepitve,
% akkor automatikusan a MiKTeX/TeX Live alap Latin Modern fontjaira valtunk.
\IfFontExistsTF{DejaVu Serif}{\setmainfont{DejaVu Serif}}{\setmainfont{Latin Modern Roman}}
\IfFontExistsTF{DejaVu Sans}{\setsansfont{DejaVu Sans}}{\setsansfont{Latin Modern Sans}}
\IfFontExistsTF{DejaVu Sans Mono}{\setmonofont{DejaVu Sans Mono}}{\setmonofont{Latin Modern Mono}}

\geometry{left=25mm,right=25mm,top=24mm,bottom=27mm}

\hypersetup{
  colorlinks=true,
  linkcolor=blue!45!black,
  urlcolor=blue!55!black,
  pdftitle={Zarovizsga tetel},
  pdfauthor={Baba Levente - tanulasi segedanyag}
}

\definecolor{accent}{HTML}{1F4E79}
\definecolor{lightaccent}{HTML}{EAF2F8}
\definecolor{warnbg}{HTML}{FFF4E5}
\definecolor{codebg}{HTML}{F7F7F7}

\tikzset{
  flow/.style={draw, rounded corners, align=center, minimum width=21mm, minimum height=8mm, fill=lightaccent},
  pred/.style={draw, diamond, aspect=2, align=center, inner sep=1pt, fill=warnbg},
  merge/.style={draw, circle, inner sep=1.2pt, fill=black!8},
  terminator/.style={draw, rounded corners=4mm, align=center, minimum width=20mm, minimum height=8mm, fill=black!5},
  arrow/.style={->, >=stealth, thick},
  zvflowstart/.style={draw, rounded rectangle, minimum width=2.3cm, minimum height=0.7cm, align=center},
  zvflowprocess/.style={draw, rectangle, minimum width=2.7cm, minimum height=0.7cm, align=center},
  zvflowdecision/.style={draw, diamond, aspect=2.2, minimum width=2.8cm, minimum height=0.8cm, align=center}
}


\titleformat{\section}{\Large\bfseries\sffamily\color{accent}}{\thesection.}{0.6em}{}
\titleformat{\subsection}{\large\bfseries\sffamily\color{accent}}{\thesubsection.}{0.6em}{}
\titleformat{\subsubsection}{\normalsize\bfseries\sffamily\color{accent}}{\thesubsubsection.}{0.6em}{}

\setlist[itemize]{topsep=3pt,itemsep=2pt,parsep=0pt,leftmargin=1.4em}
\setlist[enumerate]{topsep=3pt,itemsep=2pt,parsep=0pt,leftmargin=1.6em}
\setlength{\parindent}{0pt}
\setlength{\parskip}{6pt}
\renewcommand{\arraystretch}{1.25}
\newcolumntype{L}[1]{>{\raggedright\arraybackslash}p{#1}}
\renewcommand{\tabularxcolumn}[1]{>{\raggedright\arraybackslash}p{#1}}

\providecommand{\ZVHEADERLEFT}{\TETELNUMBER. tétel}
\providecommand{\ZVHEADERRIGHT}{\TETELSHORTTITLE}

\pagestyle{fancy}
\fancyhf{}
\lhead{\ZVHEADERLEFT}
\rhead{\ZVHEADERRIGHT}
\cfoot{\thepage}

\lstdefinelanguage{pseudocode}{
  morekeywords={IF,THEN,ELSE,WHILE,DO,FOR,TO,DOWNTO,REPEAT,UNTIL,RETURN,CALL,INPUT,OUTPUT,AND,OR,NOT,TRUE,FALSE},
  sensitive=true,
  morecomment=[l]{//}
}

\lstdefinestyle{code}{
  basicstyle=\ttfamily\small,
  backgroundcolor=\color{codebg},
  frame=single,
  rulecolor=\color{black!15},
  columns=fullflexible,
  keepspaces=true,
  breaklines=true,
  tabsize=2,
  showstringspaces=false,
  numbers=left,
  numberstyle=\tiny\color{black!55},
  numbersep=7pt,
  xleftmargin=2.2em,
  framexleftmargin=1.7em,
  captionpos=b
}

\lstdefinestyle{pseudo}{
  style=code,
  language=pseudocode,
  keywordstyle=\bfseries\color{accent!85!black},
  commentstyle=\itshape\color{black!55}
}

\lstdefinestyle{csource}{
  style=code,
  language=C,
  keywordstyle=\bfseries\color{accent!85!black},
  commentstyle=\itshape\color{black!55},
  stringstyle=\color{orange!70!black}
}

\lstset{style=code}
\renewcommand{\lstlistingname}{Kódrészlet}
\renewcommand{\lstlistlistingname}{Kódrészletek}
\newcommand{\zvcode}[2][]{\lstinputlisting[style=code,#1]{#2}}
\newcommand{\zvpseudocode}[2][]{\lstinputlisting[style=pseudo,#1]{#2}}
\newcommand{\zvccode}[2][]{\lstinputlisting[style=csource,#1]{#2}}

\newtcolorbox{summarybox}{
  enhanced,
  colback=lightaccent,
  colframe=accent!70!black,
  arc=2mm,
  boxrule=0.5pt,
  left=2mm,right=2mm,top=1mm,bottom=1mm
}

\newtcolorbox{warningbox}{
  enhanced,
  colback=warnbg,
  colframe=orange!70!black,
  arc=2mm,
  boxrule=0.5pt,
  left=2mm,right=2mm,top=1mm,bottom=1mm
}

\newcommand{\adt}[2]{\ensuremath{#1=(#2)}}
\newcommand{\set}[1]{\ensuremath{\{#1\}}}
\newcommand{\code}[1]{\texttt{#1}}
\newcommand{\base}[2]{\ensuremath{#1_{(#2)}}}
\newcommand{\addr}{\ensuremath{@}}


\providecommand{\TETELKIIRAS}{\TETELTITLE. \TETELTOPIC}
\providecommand{\ZVAUTHOR}{Baba Levente}
\providecommand{\ZVYEAR}{2026}
\providecommand{\ZVAID}{ChatGPT segítségével}

\newcommand{\makezvtitle}{%
\begin{titlepage}
  \centering
  \vspace*{18mm}
  {\Large\sffamily Programtervező Informatikus BSc\par}
  \vspace{1mm}
  {\large\sffamily \ZVYEAR-os záróvizsga tételek\par}
  \vspace{15mm}
  {\Huge\bfseries\sffamily \TETELNUMBER. tétel\par}
  \vspace{5mm}
  {\LARGE\bfseries\sffamily \TETELTITLE\par}
  \vspace{10mm}
  \begin{summarybox}
    \textbf{Tételkiírás:} \TETELKIIRAS
  \end{summarybox}
  \vfill
  {\large Készítette: \textbf{\ZVAUTHOR}\par}
  \vspace{2mm}
  {\large \ZVAID\par}
\end{titlepage}%
}

\newcommand{\zvfrontmatter}{%
  \sloppy
  \makezvtitle
  \tableofcontents
  \newpage
}

\begin{document}
\newgeometry{left=13mm,right=13mm,top=13mm,bottom=13mm}
\pagestyle{empty}
\setlength{\tabcolsep}{2.6pt}
\renewcommand{\arraystretch}{1.08}
\begin{center}
  {\Huge\bfseries Záróvizsga tételek\par}
  \vspace{2mm}
  {\Large\itshape Programtervező informatikus BSc\par}
\end{center}
\vspace{10mm}
{\fontsize{10.5}{12.2}\selectfont
\tetelsorpageone
}
\newpage
{\fontsize{10.5}{12.2}\selectfont
\tetelsorpagetwo
}
\end{document}


\newcommand{\ZVHEADERLEFT}{}
\newcommand{\ZVHEADERRIGHT}{}
\usepackage{fontspec}
\usepackage{geometry}
\usepackage{microtype}
\usepackage{xcolor}
\usepackage{amsmath,amssymb,mathtools}
\usepackage{booktabs,longtable,tabularx,array,multirow}
\usepackage{enumitem}
\usepackage{hyperref}
\usepackage{fancyhdr}
\usepackage{titlesec}
\usepackage{listings}
\usepackage[most]{tcolorbox}
\usepackage{tikz}
\usetikzlibrary{positioning,shapes.geometric,shapes.misc,arrows.meta,decorations.pathreplacing}
\usepackage{caption}

\setlength{\headheight}{16pt}
\emergencystretch=2em
\renewcommand{\contentsname}{Tartalom}

% Fontbeallitasok: ha a DejaVu fontok nincsenek telepitve,
% akkor automatikusan a MiKTeX/TeX Live alap Latin Modern fontjaira valtunk.
\IfFontExistsTF{DejaVu Serif}{\setmainfont{DejaVu Serif}}{\setmainfont{Latin Modern Roman}}
\IfFontExistsTF{DejaVu Sans}{\setsansfont{DejaVu Sans}}{\setsansfont{Latin Modern Sans}}
\IfFontExistsTF{DejaVu Sans Mono}{\setmonofont{DejaVu Sans Mono}}{\setmonofont{Latin Modern Mono}}

\geometry{left=25mm,right=25mm,top=24mm,bottom=27mm}

\hypersetup{
  colorlinks=true,
  linkcolor=blue!45!black,
  urlcolor=blue!55!black,
  pdftitle={Zarovizsga tetel},
  pdfauthor={Baba Levente - tanulasi segedanyag}
}

\definecolor{accent}{HTML}{1F4E79}
\definecolor{lightaccent}{HTML}{EAF2F8}
\definecolor{warnbg}{HTML}{FFF4E5}
\definecolor{codebg}{HTML}{F7F7F7}

\tikzset{
  flow/.style={draw, rounded corners, align=center, minimum width=21mm, minimum height=8mm, fill=lightaccent},
  pred/.style={draw, diamond, aspect=2, align=center, inner sep=1pt, fill=warnbg},
  merge/.style={draw, circle, inner sep=1.2pt, fill=black!8},
  terminator/.style={draw, rounded corners=4mm, align=center, minimum width=20mm, minimum height=8mm, fill=black!5},
  arrow/.style={->, >=stealth, thick},
  zvflowstart/.style={draw, rounded rectangle, minimum width=2.3cm, minimum height=0.7cm, align=center},
  zvflowprocess/.style={draw, rectangle, minimum width=2.7cm, minimum height=0.7cm, align=center},
  zvflowdecision/.style={draw, diamond, aspect=2.2, minimum width=2.8cm, minimum height=0.8cm, align=center}
}


\titleformat{\section}{\Large\bfseries\sffamily\color{accent}}{\thesection.}{0.6em}{}
\titleformat{\subsection}{\large\bfseries\sffamily\color{accent}}{\thesubsection.}{0.6em}{}
\titleformat{\subsubsection}{\normalsize\bfseries\sffamily\color{accent}}{\thesubsubsection.}{0.6em}{}

\setlist[itemize]{topsep=3pt,itemsep=2pt,parsep=0pt,leftmargin=1.4em}
\setlist[enumerate]{topsep=3pt,itemsep=2pt,parsep=0pt,leftmargin=1.6em}
\setlength{\parindent}{0pt}
\setlength{\parskip}{6pt}
\renewcommand{\arraystretch}{1.25}
\newcolumntype{L}[1]{>{\raggedright\arraybackslash}p{#1}}
\renewcommand{\tabularxcolumn}[1]{>{\raggedright\arraybackslash}p{#1}}

\providecommand{\ZVHEADERLEFT}{\TETELNUMBER. tétel}
\providecommand{\ZVHEADERRIGHT}{\TETELSHORTTITLE}

\pagestyle{fancy}
\fancyhf{}
\lhead{\ZVHEADERLEFT}
\rhead{\ZVHEADERRIGHT}
\cfoot{\thepage}

\lstdefinelanguage{pseudocode}{
  morekeywords={IF,THEN,ELSE,WHILE,DO,FOR,TO,DOWNTO,REPEAT,UNTIL,RETURN,CALL,INPUT,OUTPUT,AND,OR,NOT,TRUE,FALSE},
  sensitive=true,
  morecomment=[l]{//}
}

\lstdefinestyle{code}{
  basicstyle=\ttfamily\small,
  backgroundcolor=\color{codebg},
  frame=single,
  rulecolor=\color{black!15},
  columns=fullflexible,
  keepspaces=true,
  breaklines=true,
  tabsize=2,
  showstringspaces=false,
  numbers=left,
  numberstyle=\tiny\color{black!55},
  numbersep=7pt,
  xleftmargin=2.2em,
  framexleftmargin=1.7em,
  captionpos=b
}

\lstdefinestyle{pseudo}{
  style=code,
  language=pseudocode,
  keywordstyle=\bfseries\color{accent!85!black},
  commentstyle=\itshape\color{black!55}
}

\lstdefinestyle{csource}{
  style=code,
  language=C,
  keywordstyle=\bfseries\color{accent!85!black},
  commentstyle=\itshape\color{black!55},
  stringstyle=\color{orange!70!black}
}

\lstset{style=code}
\renewcommand{\lstlistingname}{Kódrészlet}
\renewcommand{\lstlistlistingname}{Kódrészletek}
\newcommand{\zvcode}[2][]{\lstinputlisting[style=code,#1]{#2}}
\newcommand{\zvpseudocode}[2][]{\lstinputlisting[style=pseudo,#1]{#2}}
\newcommand{\zvccode}[2][]{\lstinputlisting[style=csource,#1]{#2}}

\newtcolorbox{summarybox}{
  enhanced,
  colback=lightaccent,
  colframe=accent!70!black,
  arc=2mm,
  boxrule=0.5pt,
  left=2mm,right=2mm,top=1mm,bottom=1mm
}

\newtcolorbox{warningbox}{
  enhanced,
  colback=warnbg,
  colframe=orange!70!black,
  arc=2mm,
  boxrule=0.5pt,
  left=2mm,right=2mm,top=1mm,bottom=1mm
}

\newcommand{\adt}[2]{\ensuremath{#1=(#2)}}
\newcommand{\set}[1]{\ensuremath{\{#1\}}}
\newcommand{\code}[1]{\texttt{#1}}
\newcommand{\base}[2]{\ensuremath{#1_{(#2)}}}
\newcommand{\addr}{\ensuremath{@}}


\providecommand{\TETELKIIRAS}{\TETELTITLE. \TETELTOPIC}
\providecommand{\ZVAUTHOR}{Baba Levente}
\providecommand{\ZVYEAR}{2026}
\providecommand{\ZVAID}{ChatGPT segítségével}

\newcommand{\makezvtitle}{%
\begin{titlepage}
  \centering
  \vspace*{18mm}
  {\Large\sffamily Programtervező Informatikus BSc\par}
  \vspace{1mm}
  {\large\sffamily \ZVYEAR-os záróvizsga tételek\par}
  \vspace{15mm}
  {\Huge\bfseries\sffamily \TETELNUMBER. tétel\par}
  \vspace{5mm}
  {\LARGE\bfseries\sffamily \TETELTITLE\par}
  \vspace{10mm}
  \begin{summarybox}
    \textbf{Tételkiírás:} \TETELKIIRAS
  \end{summarybox}
  \vfill
  {\large Készítette: \textbf{\ZVAUTHOR}\par}
  \vspace{2mm}
  {\large \ZVAID\par}
\end{titlepage}%
}

\newcommand{\zvfrontmatter}{%
  \sloppy
  \makezvtitle
  \tableofcontents
  \newpage
}


\hypersetup{
  pdftitle={Programtervező Informatikus BSc záróvizsga tételek - teljes tételkidolgozás},
  pdfauthor={Baba Levente - ChatGPT segítségével készült tanulási segédanyag}
}

\newcommand{\ZVFullDocTitle}{Programtervező Informatikus BSc záróvizsga tételek}
\newcommand{\ZVFullDocSubtitle}{Teljes tételkidolgozás}
\newcommand{\ZVFullDocDate}{2026. június 8.}
\newcommand{\ZVCurrentHeader}{}
\newcommand{\zvsetheader}[1]{\gdef\ZVCurrentHeader{#1}}

\fancyhf{}
\fancyhead[L]{\ZVCurrentHeader}
\fancyhead[R]{Programtervező Informatikus BSc}
\fancyfoot[C]{\thepage}
\setcounter{tocdepth}{4}
\setcounter{secnumdepth}{4}
\titleformat{\paragraph}{\normalsize\bfseries\sffamily\color{accent}}{\theparagraph.}{0.6em}{}
\titlespacing*{\paragraph}{0pt}{1.0ex plus .2ex}{0.6ex}
\makeatletter
\renewcommand*{\l@section}{\@dottedtocline{1}{0em}{3.6em}}
\renewcommand*{\l@subsection}{\@dottedtocline{2}{1.5em}{4.2em}}
\renewcommand*{\l@subsubsection}{\@dottedtocline{3}{3.8em}{5.0em}}
\renewcommand*{\l@paragraph}{\@dottedtocline{4}{6.3em}{5.8em}}
\makeatother

\newcommand{\ZVFullCover}{%
\begin{titlepage}
  \thispagestyle{empty}
  \centering
  \vspace*{18mm}
  {\Large\sffamily Miskolci Egyetem\par}
  \vspace{2mm}
  {\large\sffamily Gépészmérnöki és Informatikai Kar\par}
  \vspace{18mm}
  {\Huge\bfseries\sffamily \ZVFullDocTitle\par}
  \vspace{4mm}
  {\LARGE\bfseries\sffamily \ZVFullDocSubtitle\par}
  \vspace{14mm}
  \begin{summarybox}
    Ez a dokumentum a Programtervező Informatikus BSc záróvizsgára való felkészüléshez készült. Egyetlen, közös oldalszámozású PDF-ben tartalmazza a teljes tételsort és a 27 kidolgozott tételt.
  \end{summarybox}
  \vfill
  \begin{tabular}{rl}
    \textbf{Készítette:} & Baba Levente \\
    \textbf{Készült:} & \ZVFullDocDate \\
    \textbf{Készült mire:} & Programtervező Informatikus BSc záróvizsgára való felkészülés \\
    \textbf{Segítség:} & ChatGPT segítségével készült tanulási segédanyag \\
    \textbf{Tartalom:} & teljes tételsor és 27 kidolgozott tétel \\
  \end{tabular}
  \vspace*{10mm}
\end{titlepage}%
}

\newcommand{\ZVFullTetelsorItem}[2]{%
  \item \hyperref[sec:zvfull-tetel-#1]{#2}%
}

\newcommand{\ZVFullTetelsor}{%
\section*{Tételsor}
\phantomsection
\addcontentsline{toc}{section}{Tételsor}
\begin{enumerate}[leftmargin=7mm,label=\arabic*.,itemsep=2pt,topsep=2pt,parsep=0pt]
  \ZVFullTetelsorItem{01}{\ZVTetelKiiras{1}}
  \ZVFullTetelsorItem{02}{\ZVTetelKiiras{2}}
  \ZVFullTetelsorItem{03}{\ZVTetelKiiras{3}}
  \ZVFullTetelsorItem{04}{\ZVTetelKiiras{4}}
  \ZVFullTetelsorItem{05}{\ZVTetelKiiras{5}}
  \ZVFullTetelsorItem{06}{\ZVTetelKiiras{6}}
  \ZVFullTetelsorItem{07}{\ZVTetelKiiras{7}}
  \ZVFullTetelsorItem{08}{\ZVTetelKiiras{8}}
  \ZVFullTetelsorItem{09}{\ZVTetelKiiras{9}}
  \ZVFullTetelsorItem{10}{\ZVTetelKiiras{10}}
  \ZVFullTetelsorItem{11}{\ZVTetelKiiras{11}}
  \ZVFullTetelsorItem{12}{\ZVTetelKiiras{12}}
  \ZVFullTetelsorItem{13}{\ZVTetelKiiras{13}}
  \ZVFullTetelsorItem{14}{\ZVTetelKiiras{14}}
  \ZVFullTetelsorItem{15}{\ZVTetelKiiras{15}}
  \ZVFullTetelsorItem{16}{\ZVTetelKiiras{16}}
  \ZVFullTetelsorItem{17}{\ZVTetelKiiras{17}}
  \ZVFullTetelsorItem{18}{\ZVTetelKiiras{18}}
  \ZVFullTetelsorItem{19}{\ZVTetelKiiras{19}}
  \ZVFullTetelsorItem{20}{\ZVTetelKiiras{20}}
  \ZVFullTetelsorItem{21}{\ZVTetelKiiras{21}}
  \ZVFullTetelsorItem{22}{\ZVTetelKiiras{22}}
  \ZVFullTetelsorItem{23}{\ZVTetelKiiras{23}}
  \ZVFullTetelsorItem{24}{\ZVTetelKiiras{24}}
  \ZVFullTetelsorItem{25}{\ZVTetelKiiras{25}}
  \ZVFullTetelsorItem{26}{\ZVTetelKiiras{26}}
  \ZVFullTetelsorItem{27}{\ZVTetelKiiras{27}}
\end{enumerate}%
}

\newcommand{\ZVFullIncludeTetel}[3]{%
  \clearpage
  \zvsetheader{Tétel #1}
  \section{tétel: #2}
  \label{sec:zvfull-tetel-#1}
  \begin{summarybox}
    \textbf{Tételkiírás:} \ZVTetelKiiras{#1}
  \end{summarybox}
  \begingroup
    \let\ZVOldSection\section
    \let\ZVOldSubsection\subsection
    \let\ZVOldSubsubsection\subsubsection
    \let\ZVOldParagraph\paragraph
    \renewcommand{\section}{\ZVOldSubsection}
    \renewcommand{\subsection}{\ZVOldSubsubsection}
    \renewcommand{\subsubsection}{\ZVOldParagraph}
    \input{#3}
  \endgroup
}

\begin{document}
\sloppy
\pagestyle{fancy}
\hypersetup{pageanchor=false}
\ZVFullCover

\clearpage
\hypersetup{pageanchor=true}
\pagenumbering{Roman}
\zvsetheader{Tételsor}
\ZVFullTetelsor

\clearpage
\zvsetheader{Tartalomjegyzék}
\pdfbookmark[1]{Tartalomjegyzék}{zvfull-toc}
{\small
\tableofcontents
}

\clearpage
\pagenumbering{arabic}
\setcounter{section}{0}
\ZVFullIncludeTetel{01}{Adatok, adattípusok és adatstruktúrák}{src/bodies/zv_tetel_01_body.tex}
\ZVFullIncludeTetel{02}{Algoritmusok, memória és elemi algoritmusok}{src/bodies/zv_tetel_02_body.tex}
\ZVFullIncludeTetel{03}{Strukturált programozás és programgráfok}{src/bodies/zv_tetel_03_body.tex}
\ZVFullIncludeTetel{04}{Számelméleti algoritmusok}{src/bodies/zv_tetel_04_body.tex}
\ZVFullIncludeTetel{05}{Rendezések és időelemzéseik}{src/bodies/zv_tetel_05_body.tex}
\ZVFullIncludeTetel{06}{Gráfalgoritmusok}{src/bodies/zv_tetel_06_body.tex}
\ZVFullIncludeTetel{07}{Párhuzamos algoritmusok alapjai}{src/bodies/zv_tetel_07_body.tex}
\ZVFullIncludeTetel{08}{Aggregáció, redukció, prefixszámítás és párhuzamos összefésülő rendezés}{src/bodies/zv_tetel_08_body.tex}
\ZVFullIncludeTetel{09}{Topológiák, kommunikációs modellek és aszinkron programozás}{src/bodies/zv_tetel_09_body.tex}
\ZVFullIncludeTetel{10}{Turing-gép, RAM-gép, Boole-függvények és logikai hálózatok}{src/bodies/zv_tetel_10_body.tex}
\ZVFullIncludeTetel{11}{Algoritmikus eldönthetőség, rekurzív nyelvek és polinomiális idő}{src/bodies/zv_tetel_11_body.tex}
\ZVFullIncludeTetel{12}{Nemdeterminisztikus Turing-gépek, NP, coNP és NP-teljesség}{src/bodies/zv_tetel_12_body.tex}
\ZVFullIncludeTetel{13}{Java alapok, objektumok, szemétgyűjtés és kivételkezelés}{src/bodies/zv_tetel_13_body.tex}
\ZVFullIncludeTetel{14}{Egységbezárás, információrejtés, osztálytagok és hozzáférés}{src/bodies/zv_tetel_14_body.tex}
\ZVFullIncludeTetel{15}{Öröklődés, polimorfizmus, absztrakt osztály és interfész}{src/bodies/zv_tetel_15_body.tex}
\ZVFullIncludeTetel{16}{Folyamatok, szálak és állapotmodellek az operációs rendszerekben}{src/bodies/zv_tetel_16_body.tex}
\ZVFullIncludeTetel{17}{Memóriamenedzsment, címleképzés és lapozós virtuális memória}{src/bodies/zv_tetel_17_body.tex}
\ZVFullIncludeTetel{18}{Fájlrendszerek megvalósítása}{src/bodies/zv_tetel_18_body.tex}
\ZVFullIncludeTetel{19}{Relációs adatmodell, ER modell és relációs algebra}{src/bodies/zv_tetel_19_body.tex}
\ZVFullIncludeTetel{20}{SQL nyelv, lekérdezések és relációs algebra SQL-ben}{src/bodies/zv_tetel_20_body.tex}
\ZVFullIncludeTetel{21}{Adatkezelés, adatbázis-architektúra, indexek és normalizálás}{src/bodies/zv_tetel_21_body.tex}
\ZVFullIncludeTetel{22}{Számítógép-hálózatok alapfogalmai és az ISO--OSI modell}{src/bodies/zv_tetel_22_body.tex}
\ZVFullIncludeTetel{23}{MAC alréteg, Ethernet és WLAN közeghozzáférés}{src/bodies/zv_tetel_23_body.tex}
\ZVFullIncludeTetel{24}{TCP/IP protokollszövet, Internet és IP-címzés}{src/bodies/zv_tetel_24_body.tex}
\ZVFullIncludeTetel{25}{Szoftvertechnológia és szoftverfolyamat}{src/bodies/zv_tetel_25_body.tex}
\ZVFullIncludeTetel{26}{Szoftverkövetelmények és vezérlési stílusok}{src/bodies/zv_tetel_26_body.tex}
\ZVFullIncludeTetel{27}{UML diagramok: use case, osztály-, szekvencia- és komponensdiagram}{src/bodies/zv_tetel_27_body.tex}

\end{document}
